

The report-resolved continuation theorem is exact for linear marked task equalities of a supplied time-homogeneous finite-return routine. It does not identify the reachable positive cone with every positive matrix on its support, nor replace a one-sided reachable cost bound by an unsupported operator inequality. The reset-instrument classification applies to failure branches which physically prepare known states and to first-success commitment. It includes arbitrary finite rank-one families, but not arbitrary quantum instruments or an arbitrary-rank projection without input-erasing failure. The zero-risk spectrum is exact for the acquired profile; its physical threshold has the same converse and implementation hypotheses as the exact transfer theorem. These distinctions separate the new structural results from the independent wider research goals.


The primitive finite-model route in Section~\ref{sec:primitive-risk} removes the reset and reference assumptions on its own declared interface. It supplies a common fallback, monotone two-sided approximations to the least feasible cost, and bounded-risk approximations at every strictly interior feasible budget with counted finite implementations. It does not replace the preparation-free all-initial-state quantifier by an arbitrary restricted preparation set, and it assumes finite models, finite globally legal commands/reports, time homogeneity and bounded positive call costs. It is not unknown-kernel statistical learning, a universal closed formula for the comparison profile, or optimal computational complexity. The six-regime diagram solves one explicit mixed-strata comparator with its original asymptotic qualification intact.


The compact-family route extends the primitive finite-model risk approximation to a continuously parameterized compact uncertainty class on a common finite globally legal physical interface. It controls model error and restores the same budget at every $N>N_{\min}$. Its algebraic specialization is a terminating approximate synthesis result for supplied exact primitive data, not learning from finite data, polynomial-time synthesis or a solution of the boundary finite-memory problem. The model stays fixed during a run. Exact geometric reference-profile realization, arbitrary unknown-model quotient construction and minimal computational resources remain separate mathematical questions. The examples showing boundary discontinuity and failure of vertex-only certification delimit the positive theorem rather than replace it.


The current optimality layer distinguishes three assertions. Boundary attainment and complete performed-law compactness use the shared finite-branching tree, positive cost floor and bounded terminal loss, with a fixed dictionary for total variation. Finite spectral witnesses additionally require the compact continuous primitive family and strict prefix feasibility. Same-budget finite-label approximation uses a real common fallback with strict cost slack. Hard optimized output width is not convexified, and several expected resources require a jointly safe fallback. The timing application assumes a stateless single-command device with no free clock; its exact attainment count is not a universal lower for devices that report time. None of these theorems supplies an arbitrary unknown-model quotient or a general polynomial-time synthesis method.
